\documentclass{article}
\usepackage[utf8]{inputenc}
\usepackage{amsmath}
\usepackage{geometry}
\geometry{margin=2cm}
\begin{document}

\section*{Questão 112 — ENEM 2024 — Ciências da Natureza}

\subsection*{Interpretação e Estratégia}

A questão solicita a identificação do grupo funcional na estrutura química da nimesulida capaz de reagir com uma base de Brönsted-Lowry, ou seja, que doa um próton (H$^{+}$). É necessário reconhecer o grupo que contém hidrogênio ácido.

\paragraph{Análise dos grupamentos presentes:}
\begin{itemize}
  \item Metila
  \item Sulfonamida
  \item Éter
  \item Fenila
  \item Nitro
\end{itemize}

\paragraph{Estratégia:}
\begin{enumerate}
  \item Entender que bases de Brönsted-Lowry capturam prótons.
  \item Procurar quais grupos possuem hidrogênio ácido que pode ser doado.
  \item Identificar que apenas o grupo sulfonamida possui hidrogênio ligado a nitrogênio com caráter ácido.
\end{enumerate}

\subsection*{Conteúdo e Mini Revisão}

\paragraph{Conceitos-chave}
\begin{itemize}
  \item \textbf{Base de Brönsted-Lowry:} espécie que aceita próton (H$^{+}$).
  \item \textbf{Ácido de Brönsted-Lowry:} espécie que doa próton (H$^{+}$).
  \item \textbf{Hidrogênio ácido:} hidrogênio que pode ser doado como próton, geralmente ligado a átomo eletronegativo.
\end{itemize}

\paragraph{Análise dos grupamentos:}
\begin{table}[h!]
\centering
\begin{tabular}{|l|c|c|}
\hline
Grupo & Possui hidrogênio ácido? & Justificativa \\
\hline
Metila & Não & Hidrogênios ligados a carbono alifático não são ácidos \\
Sulfonamida & Sim & Hidrogênio ligado a nitrogênio, eletropositivo e ácido \\
Éter & Não & Oxigênio ligado a carbono, sem hidrogênio ácido \\
Fenila & Não & Hidrogênios aromáticos não são ácidos \\
Nitro & Não & Não possui hidrogênio \\
\hline
\end{tabular}
\caption{Grupos funcionais da nimesulida e presença de hidrogênio ácido}
\end{table}

\subsection*{Resolução Final e Pulo do Gato}

\begin{enumerate}
  \item A nimesulida contém o grupo sulfonamida com hidrogênio ácido disponível para doação.
  \item Bases de Brönsted-Lowry capturam esse próton, formando espécie carregada.
  \item Esta reação aumenta a solubilidade da nimesulida em água.
  \item Portanto, a alternativa correta é a \textbf{A}, que indica a sulfonamida.
\end{enumerate}

\paragraph{Pulo do gato:}
Sempre procure grupos funcionais que possuem hidrogênios ácidos para responder questões de reatividade ácido-base com bases de Brönsted-Lowry.

\subsection*{Armadilhas}

\begin{itemize}
  \item Confundir grupo metila com ácido por sua visibilidade na molécula.
  \item Assumir que oxigênio implica hidrogênio ácido (éter ou nitro).
  \item Considerar hidrogênios em anel fenila ácidos.
  \item Esquecer que reações ácido-base no contexto Brönsted-Lowry envolvem transferência de próton.
\end{itemize}

\subsection*{Código da Questão}
\texttt{CN\_QUI\_ANA\_001}

\vfill
\begin{center}
\textit{Questão adaptada de:}
\newline
GONÇALVES, A. A. et al. Contextualizando reações ácido-base de acordo com a teoria protônica de Brönsted-Lowry usando comprimidos de propranolol e nimesulida. \textbf{Química Nova}, n. 3, 2013.
\end{center}


\end{document}